% Source: https://artofproblemsolving.com/community/c3h2356597_i_can_see_you_math
% Stable identity: aops-topic-2356597-q01
\textbf{(a)} Let $K$ be the set of values of $k$ allowed by the rules; by rule (1), $0\notin K$. Suppose $K$ is not empty and let $m$ be its smallest element, so $m\ge1$, and suppose the actual number of mathematicians is $k=m$.

Since $m\ge1$, it can happen that the contestant (who chooses blindly) leaves a mathematician for last. Then the $n-1$ eliminated people include exactly $m-1$ mathematicians. Seeing this, the contestant knows that $k$ is $m-1$ (if the last person is not a mathematician) or $m$ (if it is). But $m-1$ is not in $K$, because $m$ is the smallest element of $K$ (and $0\notin K$). So he is certain that $k=m$ and that the last person is a mathematician, which breaks rule (2). This contradiction shows that $K$ is empty: $\boxed{\text{there are no such pairs }(n,k)}$; the game can never be played.

(The contestant's intelligence is essential. If he used only rule (1), the pairs with $2\le k\le n$ would satisfy rule (2): after $k-1$ or $k$ eliminated mathematicians he could not tell $k-1$ from $k$, or $k$ from $k+1$. Knowing rule (2) as well, he rules out $k=1$, then $k=2$, and so on, exactly as in the argument above.)

\textbf{(b)} Suppose $D$ is not empty, and let $d$ be the latest day in $D$. If the quiz is given on day $d$, then on the morning of day $d$ there has been no quiz yet, and the students know that the quiz is on a day of $D$ that is not earlier than $d$; the only such day is $d$. So they can deduce that the quiz is that day, and $d$ does not belong to $D$, a contradiction. Hence $D$ is empty. $\blacksquare$

(Part (b) is the famous ``surprise examination'' paradox: under this model of what the students know, the announcement cannot be fulfilled, although in practice a quiz on, say, Wednesday would surprise them.)
